\section{Introduction générale}
Dans les organisations modernes, la sécurisation des réseaux informatiques est devenue un enjeu majeur face à l'augmentation des cybermenaces. Les systèmes de détection d'intrusion (IDS) permettent d'identifier les activités suspectes et les tentatives d'attaque afin de renforcer la protection des systèmes d'information.

Cependant, la diversité des solutions IDS disponibles impose aux organisations de choisir des outils adaptés à leurs besoins et à leur environnement. Ce projet a pour objectif d'évaluer et de déployer deux solutions IDS largement utilisées, Snort et Suricata, afin d'analyser leur efficacité et leurs performances à travers différents scénarios d'attaque simulés. L'étude porte notamment sur leur configuration, l'analyse des alertes générées et l'intégration d'un système d'alerte automatisé.

Le rapport présente d'abord les principales solutions IDS existantes, puis décrit les étapes de déploiement et les tests réalisés, avant de proposer une analyse comparative des résultats obtenus et des perspectives d'amélioration.
